\documentclass[]{article}

%%%%%%%%%%%%% Package imports %%%%%%%%%%%%%
\usepackage[T1]{fontenc}
\usepackage{titlesec}
\usepackage[backend=biber, style=numeric]{biblatex}
\usepackage{multicol}
\usepackage{easyReview}
\usepackage{matlab-prettifier}
\usepackage[letterpaper, margin=1in]{geometry}
\usepackage{amsmath}
\usepackage{amssymb}
\usepackage{graphicx}
\usepackage{longtable}
\usepackage{booktabs}
\usepackage{bm}
\usepackage{bbm}
\usepackage{hyperref}

%%%%%%%%%%%%% Macro definitions %%%%%%%%%%%%%

\newcommand{\pd}[2]{\frac{\partial\;#1}{\partial\;#2}}
\newcommand{\pddown}[2]{\frac{\partial}{\partial\;#2} \left[ #1 \right] }

\newcommand{\norm}[1]{\left| \left| #1 \right| \right|_{2}}

\newcommand{\x}{\bm{x}}
\newcommand{\xh}{\hat{\bm{x}}}
\newcommand{\z}{\bm{z}}
\newcommand{\Pm}{\mathbf{P}}
\newcommand{\Fm}{\mathbf{F}}
\newcommand{\Gm}{\mathbf{G}}
\newcommand{\Hm}{\mathbf{H}}
\newcommand{\Km}{\mathbf{K}}
\newcommand{\Qm}{\mathbf{Q}}
\newcommand{\Rm}{\mathbf{R}}
\newcommand{\Sm}{\mathbf{S}}
\newcommand{\Phim}{\bm{\Phi}}
\newcommand{\Id}{\mathbf{I}}

%%%%%%%%%%%%% Title card content %%%%%%%%%%%%%
\title{Extended Kalman Filter Equations with Process Noise}
\author{}
\date{}

%%%%%%%%%%%%% Begin document %%%%%%%%%%%%%
\begin{document}
\maketitle

\section*{Model}
\begin{align}
  \dot{\x}(t) &= \bm{f}\big(\x(t),\bm{u}(t),t\big) + \Gm(t)\,\bm{w}(t),
  & \mathbb{E}\!\left[\bm{w}(t)\bm{w}^\top(\tau)\right] &= \Qm(t)\,\delta(t-\tau), \\
  \z_k &= \bm{h}\big(\x(t_k),t_k\big) + \bm{v}_k,
  & \mathbb{E}\!\left[\bm{v}_k\bm{v}_j^\top\right] &= \Rm_k\,\delta_{kj},
\end{align}
with $\bm{w}$ and $\bm{v}_k$ zero-mean, mutually uncorrelated, and
uncorrelated with $\x(t_0)\sim\mathcal{N}(\xh_0,\Pm_0)$, $\Pm_0\succ 0$.

Jacobians:
\begin{equation}
  \Fm(t) = \left.\frac{\partial \bm{f}}{\partial \x}\right|_{\xh(t)}, \qquad
  \Hm_k  = \left.\frac{\partial \bm{h}}{\partial \x}\right|_{\xh_k^-}.
\end{equation}

\section*{Propagation (continuous--discrete)}
For $t\in[t_{k-1},t_k]$, starting from $\xh_{k-1}^+$ and $\Pm_{k-1}^+$:
\begin{align}
  \dot{\xh}(t) &= \bm{f}\big(\xh(t),\bm{u}(t),t\big), \\
  \dot{\Pm}(t) &= \Fm(t)\Pm(t) + \Pm(t)\Fm^\top(t) + \Gm(t)\Qm(t)\Gm^\top(t),
  \label{eq:lyapunov}
\end{align}
giving $\xh_k^- = \xh(t_k)$ and $\Pm_k^- = \Pm(t_k)$.
Equation~\eqref{eq:lyapunov} is the (linear) Lyapunov differential equation.

Equivalently, using the state transition matrix
$\dot{\Phim}(t,t_{k-1}) = \Fm(t)\Phim(t,t_{k-1})$, $\Phim(t_{k-1},t_{k-1})=\Id$:
\begin{align}
  \Pm_k^- &= \Phim_k \Pm_{k-1}^+ \Phim_k^\top + \Qm_k, \\
  \Qm_k   &= \int_{t_{k-1}}^{t_k}
             \Phim(t_k,\tau)\,\Gm(\tau)\Qm(\tau)\Gm^\top(\tau)\,\Phim^\top(t_k,\tau)\,d\tau,
\end{align}
where $\Phim_k \equiv \Phim(t_k,t_{k-1})$.

\section*{Continuous measurements: matrix Riccati ODE (Kalman--Bucy)}
If measurements are available continuously,
\begin{equation}
  \z(t) = \bm{h}\big(\x(t),t\big) + \bm{v}(t), \qquad
  \mathbb{E}\!\left[\bm{v}(t)\bm{v}^\top(\tau)\right] = \Rm(t)\,\delta(t-\tau),
  \quad \Rm(t)\succ 0,
\end{equation}
with $\bm{v}$ uncorrelated with $\bm{w}$, the extended Kalman--Bucy filter is
\begin{align}
  \dot{\xh}(t) &= \bm{f}\big(\xh(t),\bm{u}(t),t\big)
                 + \Km(t)\Big[\z(t) - \bm{h}\big(\xh(t),t\big)\Big], \\
  \Km(t) &= \Pm(t)\Hm^\top(t)\Rm^{-1}(t), \\
  \dot{\Pm}(t) &= \Fm(t)\Pm(t) + \Pm(t)\Fm^\top(t) + \Gm(t)\Qm(t)\Gm^\top(t)
                 - \Pm(t)\Hm^\top(t)\Rm^{-1}(t)\Hm(t)\Pm(t),
  \label{eq:riccati}
\end{align}
with $\Hm(t) = \partial\bm{h}/\partial\x\,\big|_{\xh(t)}$ and
$\Pm(t_0) = \Pm_0 \succ 0$.

Equation~\eqref{eq:riccati} is the matrix Riccati differential equation.
The terms $\Fm\Pm + \Pm\Fm^\top$ transport uncertainty through the dynamics,
$\Gm\Qm\Gm^\top$ injects uncertainty from process noise, and
$-\Pm\Hm^\top\Rm^{-1}\Hm\Pm$ removes uncertainty through measurement information.

\paragraph{Relation to the continuous--discrete filter.}
Between discrete measurement times no information is received, so the
quadratic term is absent and the Riccati equation reduces to the Lyapunov
equation~\eqref{eq:lyapunov} on $t\in(t_{k-1},t_k)$. The information removal
then occurs impulsively at each $t_k$ through the discrete (Joseph-form)
update below. In the limit of densely sampled measurements with
$\Rm_k = \Rm(t_k)/\Delta t$, the discrete filter recovers~\eqref{eq:riccati}.

\paragraph{Positive definiteness.}
For $\Pm_0\succ 0$ and $\Rm(t)\succ 0$, the solution of~\eqref{eq:riccati}
remains symmetric positive definite for all $t\ge t_0$ in exact arithmetic.
Numerically, integrate the $n(n+1)/2$ independent entries (or a square-root
factor) to preserve symmetry.

\section*{Measurement update (Joseph form)}
\begin{align}
  \tilde{\z}_k &= \z_k - \bm{h}\big(\xh_k^-,t_k\big)
     && \text{(innovation)} \\
  \Sm_k &= \Hm_k\Pm_k^-\Hm_k^\top + \Rm_k
     && \text{(innovation covariance)} \\
  \Km_k &= \Pm_k^-\Hm_k^\top\Sm_k^{-1}
     && \text{(gain)} \\
  \xh_k^+ &= \xh_k^- + \Km_k\tilde{\z}_k
     && \text{(state update)} \\
  \Pm_k^+ &= \big(\Id-\Km_k\Hm_k\big)\Pm_k^-\big(\Id-\Km_k\Hm_k\big)^\top
             + \Km_k\Rm_k\Km_k^\top
     && \text{(Joseph form)}
\end{align}

For the optimal gain this reduces analytically to
\begin{equation}
  \Pm_k^+ = \big(\Id-\Km_k\Hm_k\big)\Pm_k^-,
\end{equation}
but the Joseph form remains symmetric positive semidefinite under round-off
and is valid for any (possibly suboptimal) $\Km_k$.

\appendix
\section{Radial Process Noise with Variable Separation}

Let $\mathbf{d}=\mathbf{p}_2-\mathbf{p}_1$ and
$\mathbf{e}=\mathbf{d}/\norm{\mathbf{d}}$. Model uncertainty in the tether
tension as a scalar, zero-mean white Gaussian input $\xi_T(t)$ with spectral
density $q_T(t)$:
\begin{equation}
  \mathbb{E}\!\left[\xi_T(t)\xi_T(\tau)\right]
  =q_T(t)\delta(t-\tau).
\end{equation}
The deterministic accelerations are first computed from the constrained
equations, including the correction $\mathbf{q}$ in equation~13 of
\texttt{derivations/lagrangian\_based.tex}. The stochastic radial
acceleration is then added to that deterministic result, rather than passed
through the Lagrange-multiplier constraint solve.

Use normalized masses $\bar m_1=1-\eta$ and $\bar m_2=\eta$, and interpret
$\xi_T$ as an equal-and-opposite perturbing force per unit total mass. Its
acceleration contributions are obtained from the normalized force balances:
\begin{align*}
  (1-\eta)\,\delta\ddot{\mathbf{p}}_1 &= \mathbf{e}\xi_T, &
  \eta\,\delta\ddot{\mathbf{p}}_2 &= -\mathbf{e}\xi_T.
\end{align*}
Dividing each equation by its normalized mass gives the acceleration
perturbations:
\begin{align*}
  \delta\ddot{\mathbf{p}}_1 &= \frac{\mathbf{e}}{1-\eta}\xi_T, &
  \delta\ddot{\mathbf{p}}_2 &= -\frac{\mathbf{e}}{\eta}\xi_T.
\end{align*}
Since $\ddot{\mathbf{d}}=\ddot{\mathbf{p}}_2-\ddot{\mathbf{p}}_1$,
subtracting the first perturbation from the second gives the relative radial
acceleration:
\begin{equation}
  \delta\ddot{\mathbf{d}}
  =-\left(\frac{1}{\eta}+\frac{1}{1-\eta}\right)\mathbf{e}\xi_T
  =-\frac{\mathbf{e}}{\eta(1-\eta)}\xi_T.
\end{equation}
Consequently, trajectories can depart from the nominal prescribed
separation; equation~13 supplies the deterministic drift, not a hard
constraint on each noisy realization.

To place these acceleration perturbations in the EKF state equation, use the
state ordering
$\x=[\mathbf{p}_1^\top,\mathbf{p}_2^\top,
\dot{\mathbf{p}}_1^\top,\dot{\mathbf{p}}_2^\top]^\top$. In the first-order
state equation, the position derivatives are the velocities, so this
acceleration noise enters only the velocity-derivative entries. The position
blocks of the diffusion vector are zero, and its velocity blocks are the
acceleration coefficients derived above:
\begin{equation}
  \Gm_T(\x) =
  \begin{bmatrix}
    \mathbf{0}_3\\
    \mathbf{0}_3\\
    \mathbf{e}/(1-\eta)\\
    -\mathbf{e}/\eta
  \end{bmatrix},
  \qquad
  \dot{\Pm} = \Fm\Pm+\Pm\Fm^\top+\Gm_T q_T \Gm_T^\top.
\end{equation}
The covariance contribution follows by substituting the scalar input
coefficient $\Gm_T$ and its spectral density $q_T$ into the general process
noise term $\Gm\Qm\Gm^\top$ in the covariance propagation equation.

This is a nominal-drift-plus-radial-diffusion approximation: equation~13's
deterministic correction is retained, but it is not re-solved to cancel the
added noise. It is therefore not the exact fixed-length constrained model;
use it when stochastic separation variation is intentionally allowed. If a
physical tether law is required, model the tether length or compliance
explicitly and derive both drift and noise from that model. Since
$\mathbf{e}$ depends on the state, the diffusion is state-dependent; an EKF
evaluates $\Gm_T$ along the nominal trajectory, yielding a local Gaussian
approximation.

\section{Transverse Process Noise}

Let
\begin{equation}
  \mathbf{d}=\mathbf{p}_2-\mathbf{p}_1,\qquad
  \mathbf{e}=\frac{\mathbf{d}}{\norm{\mathbf{d}}},\qquad
  \mathbf{P}_\perp=\mathbf{I}_3-\mathbf{e}\mathbf{e}^\top.
\end{equation}
The matrix $\mathbf{P}_\perp$ projects onto the plane perpendicular to the
tether direction. Model an uncertain transverse force per unit total mass
using a zero-mean white Gaussian vector input $\bm{\xi}_\perp(t)$:
\begin{equation}
  \mathbb{E}\!\left[
    \bm{\xi}_\perp(t)\bm{\xi}_\perp^\top(\tau)
  \right]
  =q_\perp(t)\mathbf{I}_3\delta(t-\tau).
\end{equation}
The equal-and-opposite normalized force perturbations are
\begin{equation}
  \delta\mathbf{f}_1(t)=\mathbf{P}_\perp\bm{\xi}_\perp(t),\qquad
  \delta\mathbf{f}_2(t)=-\mathbf{P}_\perp\bm{\xi}_\perp(t).
\end{equation}
Although the input is written as a three-vector, projection removes its
component along $\mathbf{e}$, leaving two independent transverse directions.

Using normalized masses $\bar m_1=1-\eta$ and $\bar m_2=\eta$, divide the
force perturbations by the respective masses to obtain
\begin{align}
  \delta\ddot{\mathbf{p}}_1
    &=\frac{\mathbf{P}_\perp\bm{\xi}_\perp}{1-\eta}, &
  \delta\ddot{\mathbf{p}}_2
    &=-\frac{\mathbf{P}_\perp\bm{\xi}_\perp}{\eta}, \\
  \delta\ddot{\mathbf{d}}
    &=-\left(\frac{1}{1-\eta}+\frac{1}{\eta}\right)
      \mathbf{P}_\perp\bm{\xi}_\perp
    =-\frac{\mathbf{P}_\perp\bm{\xi}_\perp}{\eta(1-\eta)}.
\end{align}
The relative acceleration is perpendicular to the current separation,
because $\mathbf{d}^\top\mathbf{P}_\perp=\mathbf{0}^\top$. Thus this noise
has no instantaneous radial-acceleration component.

For the state ordering
$\x=[\mathbf{p}_1^\top,\mathbf{p}_2^\top,
\dot{\mathbf{p}}_1^\top,\dot{\mathbf{p}}_2^\top]^\top$, define
\begin{equation}
  \Gm_\perp(\x)=
  \begin{bmatrix}
    \mathbf{0}_{3\times3}\\
    \mathbf{0}_{3\times3}\\
    \mathbf{P}_\perp/(1-\eta)\\
    -\mathbf{P}_\perp/\eta
  \end{bmatrix},\qquad
  \dot{\Pm}=\Fm\Pm+\Pm\Fm^\top
    +\Gm_\perp(q_\perp\mathbf{I}_3)\Gm_\perp^\top.
\end{equation}
For this purely transverse input, passing the unconstrained acceleration
noise through the acceleration map in equation~13 gives the same result:
the map corrects only the relative component parallel to $\mathbf{d}$, and
$\mathbf{P}_\perp$ has no such component. With an initially constraint-
consistent state and deterministic drift that enforces the acceleration-
level constraint, this continuous-time diffusion does not inject noise into
the fixed-length or radial-velocity constraints. A finite-step numerical
integrator may nevertheless accumulate constraint error, which should be
monitored (and, if necessary, controlled with a constraint-consistent
integrator). Since $\mathbf{P}_\perp$ depends on the state, an EKF evaluates
$\Gm_\perp$ along the nominal trajectory.

\section{Sensors}\label{sec:sensors}

Table~\ref{tab:sensor_list} lists candidate sensors; the model listing below
defines their idealized measurement equations. The current implementation in
\texttt{coa.meas.measurementModels} covers sensors 1--19 in eight selectable
groups. Sensors 20--22 are accelerometer candidates only and are not currently
implemented.

{\footnotesize
\setlength{\tabcolsep}{3pt}
\renewcommand{\arraystretch}{1.15}
\begin{longtable}{@{}r p{0.38\linewidth} p{0.20\linewidth} p{0.16\linewidth} p{0.12\linewidth}@{}}
\caption{Candidate sensors.}\label{tab:sensor_list}\\
\toprule
\# & Sensor & Sensor model & Jacobian & Cadence \\
\midrule
\endfirsthead
\toprule
\# & Sensor & Sensor model & Jacobian & Cadence \\
\midrule
\endhead
\bottomrule
\endlastfoot
1 & Ground optical angles, $\mathbf{r}_1$ & $\begin{aligned}\alpha &= h_1(\mathbf{r}_1) \\ \delta &= h_2(\mathbf{r}_1) \end{aligned}$ & & \\
\midrule
2 & Ground optical angles, $\mathbf{r}_2$ & $\begin{aligned}\alpha &= h_1(\mathbf{r}_2) \\ \delta &= h_2(\mathbf{r}_2) \end{aligned}$ & & \\
\midrule
3 & Ground optical angles, $\mathbf{r}_3$ & $\begin{aligned}\alpha &= h_1(\mathbf{r}_3) \\ \delta &= h_2(\mathbf{r}_3) \end{aligned}$ & & \\
\midrule
4 & Optical relative-geometry angles, $\mathbf{r}_2 - \mathbf{r}_3$ & $\begin{aligned}\alpha &= h_1(\mathbf{r}_2 - \mathbf{r}_3) \\ \delta &= h_2(\mathbf{r}_2 - \mathbf{r}_3) \end{aligned}$ & & \\
\midrule
5 & Star trackers, tether orientation, $\mathbf{r}_2 - \mathbf{r}_1$ & $\begin{aligned}\alpha &= h_1(\mathbf{r}_2 - \mathbf{r}_1) \\ \delta &= h_2(\mathbf{r}_2 - \mathbf{r}_1) \end{aligned}$ & & \\
\midrule
6 & Optical angular size of tether in payload sky, $\mathbf{r}_1, \mathbf{r}_2, \mathbf{r}_3$ & $\theta = h_3(\mathbf{r}_1, \mathbf{r}_2, \mathbf{r}_3)$ & & \\
\midrule
7 & DSN ground radar range, $\mathbf{r}_1$ & $r = h_4(\mathbf{r}_1)$ & & \\
\midrule
8 & DSN ground radar range, $\mathbf{r}_2$ & $r = h_4(\mathbf{r}_2)$ & & \\
\midrule
9 & DSN ground radar range, $\mathbf{r}_3$ & $r = h_4(\mathbf{r}_3)$ & & \\
\midrule
10 & DSN ground radar range rate, $\mathbf{r}_1$ & $\dot{r} = h_5(\mathbf{r}_1, \dot{\mathbf{r}}_1)$ & & \\
\midrule
11 & DSN ground radar range rate, $\mathbf{r}_2$ & $\dot{r} = h_5(\mathbf{r}_2, \dot{\mathbf{r}}_2)$ & & \\
\midrule
12 & DSN ground radar range rate, $\mathbf{r}_3$ & $\dot{r} = h_5(\mathbf{r}_3, \dot{\mathbf{r}}_3)$ & & \\
\midrule
13 & Dedicated RF ranging crosslink, tether length, $\mathbf{r}_2 - \mathbf{r}_1$ & $ r = h_4(\mathbf{r}_2 - \mathbf{r}_1)$ & & \\
\midrule
14 & Doppler from RF ranging crosslink, $\dot{L}$, $\mathbf{r}_2 - \mathbf{r}_1$ & $\dot{r} = h_5(\mathbf{r}_2 - \mathbf{r}_1, \dot{\mathbf{r}}_2 - \dot{\mathbf{r}}_1)$ & & \\
\midrule
15 & Dedicated RF ranging crosslink, payload relative position, $\mathbf{r}_3 - \mathbf{r}_2$ & $ r = h_4(\mathbf{r}_3 - \mathbf{r}_2)$ & & \\
\midrule
16 & Doppler from RF ranging crosslink, payload relative in-track speed, $\mathbf{r}_3 - \mathbf{r}_2$ & $\dot{r} = h_5(\mathbf{r}_3 - \mathbf{r}_2, \dot{\mathbf{r}}_3 - \dot{\mathbf{r}}_2)$ & & \\
\midrule
17 & Lunar PN\&T/GNSS-like position fix, $\mathbf{r}_1$ & $\mathbf{r}_1$ & & \\
\midrule
18 & Lunar PN\&T/GNSS-like position fix, $\mathbf{r}_2$ & $\mathbf{r}_2$ & & \\
\midrule
19 & Lunar PN\&T/GNSS-like position fix, $\mathbf{r}_3$ & $\mathbf{r}_3$ & & \\
\midrule
20 & Candidate accelerometer on $m_1$ (not implemented) & Noise plus bias & & \\
\midrule
21 & Candidate accelerometer on $m_2$ (not implemented) & Noise plus bias & & \\
\midrule
22 & Candidate accelerometer on $m_3$ (not implemented) & Noise plus bias & & \\
\end{longtable}
}

\subsection{Model listing}

Angles measurement, input position vector $\boldsymbol{\rho}$:

\begin{align}
  \alpha &= h_1(\boldsymbol{\rho})
          = \operatorname{atan2}(\rho_y,\rho_x), \\
  \delta &= h_2(\boldsymbol{\rho})
          = \operatorname{atan2}\!\left(\rho_z,
              \sqrt{\rho_x^2+\rho_y^2}\right).
  \label{eq:measurement-angles}
\end{align}
The declination expression is equivalent to
$\arcsin(\rho_z/\|\boldsymbol{\rho}\|)$. Right ascension is quadrant-aware;
the value and its Jacobian are undefined on the $z$ axis
($\rho_x=\rho_y=0$).

The Jacobian returned by \texttt{coa.meas.angles} is
\begin{equation}
  \pd{[\alpha,\delta]^\top}{\boldsymbol{\rho}} =
  \begin{bmatrix}
    -\rho_y/q & \rho_x/q & 0 \\
    -\rho_x\rho_z/(r^2\sqrt{q}) &
    -\rho_y\rho_z/(r^2\sqrt{q}) & \sqrt{q}/r^2
  \end{bmatrix},\quad
  q=\rho_x^2+\rho_y^2,\quad r=\|\boldsymbol{\rho}\|.
\end{equation}

Angular size measurement, input position vectors $\boldsymbol{\rho}_1$ and
$\boldsymbol{\rho}_2$:

\begin{equation}
  \theta = h_3(\boldsymbol{\rho}_1,\boldsymbol{\rho}_2)
  = \arccos(\mathbf{u}_1^\top\mathbf{u}_2),\qquad
  \mathbf{u}_i=\frac{\boldsymbol{\rho}_i}{\|\boldsymbol{\rho}_i\|}.
  \label{eq:measurement-angular-size}
\end{equation}
The cosine is clamped to $[-1,1]$ against floating-point roundoff. The
measurement value is defined for parallel and antiparallel nonzero vectors,
but the Jacobians are rejected when $\sin\theta<10^{-12}$. For other inputs,
\begin{align}
  \pd{\theta}{\boldsymbol{\rho}_1}
    &= -\frac{(\mathbf{u}_2-c\mathbf{u}_1)^\top}
             {\|\boldsymbol{\rho}_1\|\sin\theta}, \\
  \pd{\theta}{\boldsymbol{\rho}_2}
    &= -\frac{(\mathbf{u}_1-c\mathbf{u}_2)^\top}
             {\|\boldsymbol{\rho}_2\|\sin\theta},\qquad
  c=\mathbf{u}_1^\top\mathbf{u}_2.
\end{align}

Range measurement, input position vector $\boldsymbol{\rho}$:

\begin{equation}
  r = h_4(\boldsymbol{\rho}) = \|\boldsymbol{\rho}\|,\qquad
  \pd{r}{\boldsymbol{\rho}} = \frac{\boldsymbol{\rho}^\top}{r}
  \quad(r>0).
  \label{eq:measurement-range}
\end{equation}
The range value is defined at $\boldsymbol{\rho}=0$, but the Jacobian is not.

Range rate measurement, input position and velocity vectors
$\boldsymbol{\rho}$ and $\dot{\boldsymbol{\rho}}$:

\begin{align}
  \dot{r} &= h_5(\boldsymbol{\rho},\dot{\boldsymbol{\rho}})
    = \frac{\boldsymbol{\rho}^\top\dot{\boldsymbol{\rho}}}{r},
    \qquad r=\|\boldsymbol{\rho}\|, \\
  \pd{\dot{r}}{\boldsymbol{\rho}}
    &= \frac{(\dot{\boldsymbol{\rho}}-\dot{r}\mathbf{u})^\top}{r},
  \qquad
  \pd{\dot{r}}{\dot{\boldsymbol{\rho}}}=\mathbf{u}^\top,
  \qquad \mathbf{u}=\frac{\boldsymbol{\rho}}{r}.
  \label{eq:measurement-range-rate}
\end{align}
Both value and Jacobians are undefined at $r=0$.

The lunar PN\&T position measurement is
\begin{equation}
  \mathbf{z}_{\mathrm{PN\&T}}=h_6(\boldsymbol{\rho})
  =\boldsymbol{\rho},\qquad
  \pd{\mathbf{z}_{\mathrm{PN\&T}}}{\boldsymbol{\rho}}=\mathbf{I}_3.
  \label{eq:measurement-pnt}
\end{equation}
It measures position only. The current model does not include an
accelerometer/IMU measurement equation or bias state for candidate sensors
20--22.

\end{document}
